\documentclass[../MAIN.tex]{subfiles}

\begin{document}
	
	% ============================================================================
	% COMPARATIVA INICIAL DE TRACKERS
	% ============================================================================
	
	\section{Comparativa inicial de trackers}
	
	% Materiales base:
	% - evaluation_results.md
	% - evaluation_results_1.md
	% - evaluation_results_2.md
	% - conversaciones iniciales con tutores
	% - analisis_bytetrack_motip.md (solo como apoyo interpretativo)
	
	\begin{chaptertemplatebox}
		\textbf{Objetivo del experimento}\\
		Comparar varios trackers representativos sobre una misma secuencia de fútbol para identificar cuál ofrecía el mejor punto de partida para el desarrollo posterior del TFG.\\[0.5em]
		
		\textbf{Material usado}\\
		Tablas de evaluación inicial sobre SNMOT-060, vídeos generados en las primeras pruebas y conversaciones con tutores sobre paradigmas de tracking tradicionales frente a métodos end-to-end.\\[0.5em]
		
		\textbf{Configuración}\\
		Secuencia SNMOT-060 de SoccerNet, aproximadamente 750 frames y 30 segundos de duración. Se evaluaron ByteTrack sin ReID, ByteTrack con ReID OSNet entrenado en SportsMOT, DeepEIoU con el mismo modelo ReID y MOTIP con pesos oficiales entrenados sobre SportsMOT. La evaluación se realizó con métricas MOT estándar y, posteriormente, con métricas ampliadas que incluían fragmentaciones y número total de IDs.\\[0.5em]
		
		\textbf{Resultados principales}\\
		ByteTrack sin ReID mostró el mejor equilibrio global en la comparación inicial, mientras que ByteTrack con ReID redujo cambios de identidad e IDs totales a costa de una ligera degradación en métricas globales. DeepEIoU presentó una fragmentación claramente superior y MOTIP mostró peor estabilidad global en este dominio.\\[0.5em]
		
		\textbf{Interpretación}\\
		La comparación inicial sugirió que el problema no era únicamente “usar o no usar ReID”, sino elegir una base modular suficientemente sólida sobre la que analizar después la asociación de identidad.\\[0.5em]
		
		\textbf{Limitaciones}\\
		Estas pruebas corresponden a una fase inicial del trabajo y todavía no incorporan ni la recalibración del evaluador ni las mejoras posteriores del pipeline. Por tanto, deben leerse como una comparación exploratoria para selección de base, no como resultados finales del TFG.\\[0.5em]
		
		\textbf{Conclusión parcial}\\
		ByteTrack fue seleccionado como tracker base del trabajo por ofrecer el mejor equilibrio entre rendimiento, interpretabilidad y capacidad de modificación experimental.
	\end{chaptertemplatebox}
	
	\subsection{Selección de trackers candidatos}
	
	La primera fase experimental del trabajo se planteó como una comparación entre varios trackers representativos de paradigmas distintos. El objetivo no era todavía optimizar un sistema final, sino determinar qué enfoque ofrecía una base más sólida para el dominio del fútbol. En las conversaciones iniciales con los tutores se planteó explícitamente comparar métodos tradicionales del paradigma \textit{tracking-by-detection}, donde el tracking y la asociación pueden modularse por separado, frente a enfoques más integrados o \textit{end-to-end}, como MOTIP, donde la identidad se aprende de forma más directa dentro del propio sistema.
	
	Siguiendo este criterio, se seleccionaron cuatro configuraciones iniciales. La primera fue ByteTrack sin ReID, utilizada como línea base del paradigma modular. La segunda fue ByteTrack con un modelo ReID OSNet entrenado en SportsMOT, con el fin de comprobar si la incorporación de apariencia mejoraba la continuidad de identidad. La tercera fue DeepEIoU, también combinado con el mismo modelo ReID, como alternativa basada en un esquema de tracking deportivo con distinta política de asociación. Finalmente, se incluyó MOTIP como representante de un enfoque más reciente basado en transformers y predicción de identidad.
	
	Esta selección resultó metodológicamente útil porque permitía responder simultáneamente a dos preguntas. Por un lado, si un tracker modular fuerte como ByteTrack era suficiente como base de trabajo. Por otro, si un método más integrado como MOTIP podía ofrecer ventajas en un escenario deportivo complejo sin necesidad de incorporar un módulo ReID explícito.
	
	La Tabla~\ref{tab:comparativa_inicial_trackers} resume los resultados cuantitativos de la comparación inicial entre los trackers evaluados. Como puede observarse, ByteTrack sin ReID ofrecía ya un equilibrio muy competitivo entre métricas globales y estabilidad de identidad, mientras que la incorporación de ReID reducía notablemente los cambios de identidad y el número total de trayectorias. En cambio, DeepEIoU y MOTIP mostraban una fragmentación claramente superior, lo que justificó descartar ambos como base principal del desarrollo posterior.
	
	\begin{table}[H]
		\centering
		\caption{Comparativa inicial de trackers sobre la secuencia SNMOT-060.}
		\label{tab:comparativa_inicial_trackers}
		\begin{tabular}{lccccc}
			\toprule
			\textbf{Método} & \textbf{MOTA} & \textbf{IDF1} & \textbf{ID Switches} & \textbf{Frags} & \textbf{Total IDs} \\
			\midrule
			ByteTrack sin ReID                & 80.7\% & 75.6\% & 41 & 119 & 62 \\
			ByteTrack + ReID (SportsMOT)      & 79.2\% & 74.4\% & 36 & 131 & 30 \\
			DeepEIoU + ReID (SportsMOT)       & 80.9\% & 75.2\% & 60 & 154 & 67 \\
			MOTIP                             & 73.1\% & 71.0\% & 51 & 350 & 61 \\
			\bottomrule
		\end{tabular}
	\end{table}
	
	Además de las configuraciones recogidas en la Tabla~\ref{tab:comparativa_inicial_trackers}, se realizaron otras pruebas preliminares adicionales durante la fase de selección inicial. Estas se incluyen en el Apéndice~\ref{ap:tablas_completas} para no sobrecargar la narrativa principal.
	
	En conjunto, la tabla muestra que ByteTrack era la base más equilibrada para continuar el trabajo. Aunque DeepEIoU alcanzaba un valor de MOTA ligeramente superior, su número de cambios de identidad y fragmentaciones era mucho mayor. Por su parte, MOTIP presentaba un comportamiento claramente más inestable en este dominio. ByteTrack con ReID ya apuntaba, además, que la señal de apariencia podía aportar valor real al reducir los ID switches y compactar el número total de trayectorias.
	
	\subsection{ByteTrack}
	
	ByteTrack se utilizó como referencia principal de la familia \textit{tracking-by-detection}. La razón de esta elección fue doble. En primer lugar, se trata de un tracker con muy buen rendimiento en la literatura reciente y con una filosofía especialmente interesante para escenarios con detecciones degradadas: en vez de descartar de manera temprana las cajas de baja confianza, intenta asociarlas también, recuperando instancias verdaderas que de otro modo fragmentarían la trayectoria. En segundo lugar, su implementación permite intervenir de forma relativamente sencilla en parámetros, lógica de asociación y acoplamiento con módulos ReID, lo que lo convertía en una base muy adecuada para un TFG centrado en análisis y modificación del pipeline.
	
	En la comparación inicial, ByteTrack sin ReID obtuvo un rendimiento cuantitativo muy competitivo: 80.7\% de MOTA, 75.6\% de IDF1 y 41 ID switches. Además, en la tabla ampliada aparecía con 119 fragmentaciones y 62 IDs totales, lo que evidenciaba que, aunque la calidad global del tracking era alta, seguía existiendo un problema apreciable de fragmentación de identidades. Esta combinación de buen rendimiento base y margen claro de mejora fue precisamente una de las razones por las que ByteTrack resultó especialmente atractivo como punto de partida.
	
	Cuando se añadió un modelo OSNet entrenado en SportsMOT, ByteTrack pasó a 79.2\% de MOTA y 74.4\% de IDF1, pero redujo los ID switches de 41 a 36 y, sobre todo, redujo drásticamente los IDs totales de 62 a 30. Esta observación fue importante desde el principio, porque sugería que la señal ReID sí estaba aportando continuidad de identidad, aunque su integración no resultaba todavía óptima y generaba un trade-off entre compactación de identidades y calidad global del seguimiento.
	
	\subsection{DeepEIoU}
	
	DeepEIoU se incluyó como tracker alternativo dentro del conjunto de métodos probados en la fase inicial. Su interés residía en comprobar si un enfoque ya orientado a escenarios deportivos podía superar a ByteTrack al combinar asociación geométrica y apariencia. Para mantener la comparabilidad, se utilizó también un modelo OSNet entrenado en SportsMOT.
	
	Los resultados cuantitativos mostraron que DeepEIoU alcanzaba 80.9\% de MOTA y 75.2\% de IDF1, es decir, un valor de MOTA ligeramente superior al de ByteTrack sin ReID. Sin embargo, esta aparente ventaja desaparecía al observar las métricas de identidad: 60 ID switches, 154 fragmentaciones y 67 IDs totales. En otras palabras, el tracker mantenía una calidad global aceptable en términos de detección y cobertura, pero lo hacía a costa de una estabilidad de identidad claramente peor.
	
	Este comportamiento fue muy relevante para la selección del sistema base. La comparación mostraba que una métrica global ligeramente mejor no implicaba necesariamente un mejor tracker para el problema concreto del TFG. En un dominio donde la continuidad de identidad es un objetivo central, el número de cambios de ID y la fragmentación pesan más que una mejora marginal en MOTA. Por ello, aunque DeepEIoU se conservó como referencia comparativa, no se consideró la mejor base para el desarrollo posterior.
	
	\subsection{MOTIP}
	
	MOTIP se incorporó a la comparativa como representante de un paradigma distinto al de los trackers modulares clásicos. A diferencia de ByteTrack o DeepEIoU, no sigue la filosofía de \textit{tracking-by-detection} con un módulo ReID acoplado, sino que plantea el seguimiento como un problema de predicción de identidad más integrado. Precisamente por ello resultaba interesante incluirlo en la fase inicial: permitía contrastar un método más reciente y cercano al paradigma \textit{end-to-end} frente a los trackers tradicionales.
	
	Sin embargo, los resultados cuantitativos iniciales fueron claramente peores que los de ByteTrack. MOTIP obtuvo 73.1\% de MOTA, 71.0\% de IDF1 y 51 ID switches. En la tabla ampliada, además, aparecía con 350 fragmentaciones y 61 IDs totales, cifras que reflejaban una inestabilidad notable de las trayectorias. Esta degradación no solo se observó en las métricas, sino también en el análisis visual posterior, donde se documentaron fenómenos de reciclaje de identidades, teletransportación por paneo de cámara y creación de IDs fantasma durante oclusiones y caídas de jugadores.
	
	Desde el punto de vista metodológico, la inclusión de MOTIP fue, aun así, muy valiosa. Permitió comprobar que el paradigma alternativo era interesante desde un punto de vista conceptual, pero que en el dominio concreto del fútbol y bajo el protocolo experimental del TFG no constituía la mejor base para construir mejoras progresivas. Además, su menor modularidad lo hacía menos conveniente para un trabajo cuyo objetivo no era solo comparar resultados, sino intervenir explícitamente en la lógica de asociación e identidad.
	
	\subsection{Resultados iniciales y selección del tracker base}
	
	La comparación inicial permite extraer una conclusión clara. Si se observa únicamente MOTA, ByteTrack sin ReID y DeepEIoU ofrecen valores similares, e incluso DeepEIoU aparece ligeramente por encima. Sin embargo, al incorporar métricas de identidad y fragmentación, la diferencia entre métodos se vuelve mucho más clara. ByteTrack sin ReID logra el mejor equilibrio entre rendimiento global y estabilidad de identidad dentro de la comparación inicial, mientras que ByteTrack con ReID muestra ya una reducción notable de ID switches e IDs totales, aunque todavía con un trade-off no resuelto en MOTA e IDF1.
	
	En cambio, DeepEIoU y MOTIP quedan por detrás cuando el criterio se desplaza desde la mera cobertura global hacia la continuidad real de identidad. DeepEIoU incrementa fuertemente los cambios de ID y la fragmentación, y MOTIP presenta un comportamiento todavía más inestable, tanto cuantitativa como visualmente. Por ello, ambos métodos se conservaron como referencia comparativa, pero no como base principal del desarrollo posterior.
	
	La decisión final de esta fase fue seleccionar ByteTrack como tracker base del TFG. Esta elección se apoyó en tres razones. En primer lugar, ofrecía el mejor equilibrio inicial entre métricas globales y continuidad de identidad. En segundo lugar, permitía experimentar de forma controlada con la integración de ReID, el ajuste de memoria temporal y la modificación de la función de asociación. Finalmente, su carácter modular lo hacía compatible con el tipo de contribución que se buscaba en este trabajo: no solo evaluar trackers existentes, sino analizar por qué fallan y proponer mecanismos concretos para mejorar su comportamiento en fútbol.
	
	\newpage
	
\end{document}